% !TeX root = ../../Book2.tex
\chapter{\texorpdfstring{\(L^2\)}{L²} 与 \texorpdfstring{\(L^p\)}{Lᵖ} Fourier 理论}
\label{b2:ch:04}
\BookChapterNumber{b2:ch:04}
% Axler2020Measure, Section11C, 印刷375--377页核对L2延拓的定义和证明路线。
% Hausdorff--Young在本章用有限简单函数和已证明三线引理完整推导；不前借第19章。
\begin{chapterabstract}
Fourier 变换不仅处理可积函数，还可以通过能量等式作用在整个平方可积空间。本章证明 Plancherel 延拓的良定义、满射性与反演，再用两端估计和复参数插值建立 Hausdorff--Young 不等式。空间伸缩与分离波包的例子显示，大于二的指数不能按同样方式延拓到共轭指数空间。最后以频率乘子说明酉性怎样简化线性算子，并引出广义变换的需要。
\end{chapterabstract}
\begin{learninggoals}
\begin{itemize}
\item\label{b2:item:04-goal-extension}从共同稠密类构造唯一的平方范数变换，说明定义与逼近列无关。
\item\label{b2:item:04-goal-unitary}用反演证明酉性，区分等距、满射及逐点积分三个不同问题。
\item\label{b2:item:04-goal-interpolation}理解两端估计怎样经参数函数产生中间指数，并核对插值的每个积分和极限。
\item\label{b2:item:04-goal-failure}用尺度及符号平均构造 \(p>2\) 的范数障碍，准确表述其否定范围。
\item\label{b2:item:04-goal-multiplier}在 \(L^2\) 上定义频率乘子，计算范数，说明推广到其他指数还需哪些估计。
\end{itemize}
\end{learninggoals}

测量一个信号时，平方积分常表示它的总能量。如果把信号分成不同频率，能量是否恰好等于频率幅值的平方积分？这个问题对不可积但平方可积的信号同样有意义；然而上一章的绝对收敛积分还不能定义它们的变换。本章从能量等式出发构造变换，并用它求解 \((a-\Delta)u=f\) 这样的方程。之后再问：当输入只有 \(p\) 次可积性时，频率幅值能落在哪个空间？

\ProseParagraphBreak

从可积函数出发得到的反演公式，怎样适用于仅有有限能量的信号？可以先选容易计算的函数逼近信号，再用能量等式控制变换的误差。这个构造将积分公式与平方范数极限联系起来。

\ProseParagraphBreak

\secref{b2:sec:0401}建立能量等式，\secref{b2:sec:0402}再用复参数把可积端点与平方端点连接起来。\secref{b2:sec:0403}说明所得结论的适用边界。Sheldon Axler 的《Measure, Integration \& Real Analysis》\cite[Section 11C, pp. 375--377]{Axler2020Measure}给出了 Plancherel 延拓的另一种讲法；更系统的 \(L^p\) Fourier 理论见 Loukas Grafakos 的《Classical Fourier Analysis》\cite{Grafakos2014Classical}。

% !TeX root = ../../Book2.tex
\section{Plancherel 定理}
\label{b2:sec:0401}

上一章把 Fourier 变换定义为绝对收敛积分，因此要求 \(f\in L^1\)。平方可积函数未必可积，但其能量仍有限。本节先在两个空间的交集上证明能量守恒，再利用完备性延拓变换。延拓所得对象是一个 \(L^2\) 等价类，不能自动解释为每个频率处的普通积分。

\subsection{\texorpdfstring{\(L^1\cap L^2\)}{L¹∩L²} 上的平方范数恒等式}
\label{b2:subsec:040101}

先在同时具有可积变换的函数上计算。

\ProseParagraphBreak

有限区间上的正交展开把平方范数写成系数平方和。在整个空间上，离散的频率指标变为连续变量，所期待的等式是两个积分相等。这里需要先处理一个实际障碍：由 \(f\in L^1\cap L^2\) 只能立即知道 \(\widehat f\) 连续有界，不能立即知道它平方可积。以下第一步先把两侧的积分合法性建立起来，第二步才让正则化消失。

\begin{lemma}\label{b2:lem:04-parseval-core}
设 \(d\) 为正整数，\(f,g\in L^1(\mathbb R^d;\mathbb C)\cap L^2(\mathbb R^d;\mathbb C)\)，且 \(\widehat f\in L^1(\mathbb R^d)\)。则
\[
 \int f(x)\cdot\overline{g(x)}\,\mathrm{d}x
 =\int\widehat f(\xi)\cdot\overline{\widehat g(\xi)}\,\mathrm{d}\xi.
\]
\end{lemma}
\begin{proof}
左侧由 Cauchy--Schwarz 不等式绝对收敛。对 \(f\) 使用上一章的反演定理，使 \(f(x)=\int\widehat f(\xi)e^{2\pi\mathrm{i}x\cdot\xi}\,\mathrm{d}\xi\) 几乎处处。将此式乘以 \(\overline g(x)\) 再积分；重积分绝对值不超过 \(\|\widehat f\|_1\cdot\|g\|_1\)，故可交换次序。内层积分为 \(\overline{\widehat g(\xi)}\)。右侧绝对收敛还可由 \(\widehat f\in L^1\)、\(\|\widehat g\|_\infty\leq\|g\|_1\) 直接确认。
\end{proof}

\begin{theorem}[Plancherel，交空间形式]\label{b2:thm:04-plancherel-core}
设 \(d\) 为正整数。若 \(f\in L^1(\mathbb R^d;\mathbb C)\cap L^2(\mathbb R^d;\mathbb C)\)，则 \(\widehat f\in L^2(\mathbb R^d;\mathbb C)\)，且 \(\|\widehat f\|_2=\|f\|_2\)。
\end{theorem}
\begin{proof}
令 \(f_\varepsilon=G_\varepsilon*f\)，其中 \(G_\varepsilon(x)=\varepsilon^{-d}e^{-\pi|x|^2/\varepsilon^2}\)，积分为 \(1\)。Young 不等式分别给出 \(\|f_\varepsilon\|_1\leq\|f\|_1\)、\(\|f_\varepsilon\|_2\leq\|f\|_2\)；\secref{b2:sec:0304}的近似恒等定理给出 \(f_\varepsilon\to f\) 于 \(L^2\)。卷积定理给出
\[
 \widehat {f_\varepsilon}(\xi)
 =e^{-\pi\varepsilon^2|\xi|^2}\cdot\widehat f(\xi).
\]
右侧是有界函数与可积 Gauss 因子的乘积，所以属于 \(L^1\)，也属于 \(L^2\)。对 \(f_\varepsilon\) 使用前一引理，得到
\[
 \|f_\varepsilon\|_2^2
 =\int e^{-2\pi\varepsilon^2|\xi|^2}|\widehat f(\xi)|^2\,\mathrm{d}\xi.
\]
取任意递减到零的 \(\varepsilon_n\)。范数的反三角不等式使左侧趋于 \(\|f\|_2^2\)；右侧的非负被积函数逐点递增到 \(|\widehat f|^2\)，由单调收敛趋于 \(\int|\widehat f|^2\)，即使预先尚不知道此积分有限也可使用。左侧极限有限，故同时得到平方可积性和等式。由于目标等式不含所选序列，它适用于原来的每个 \(f\)。
\end{proof}

这称为\term[plancherel-dingli]{Plancherel 定理}[Plancherel's theorem]。Gauss 因子先使变换可积，反演与积分换序便可在可积控制下进行，最后再取正则化参数的极限。

\subsection{Plancherel 延拓}
\label{b2:subsec:040102}

\begin{theorem}\label{b2:thm:04-plancherel-extension}
设 \(d\) 为正整数。存在唯一复线性等距映射 \(\mathcal F_2:L^2(\mathbb R^d;\mathbb C)\to L^2(\mathbb R^d;\mathbb C)\)，在 \(L^1(\mathbb R^d;\mathbb C)\cap L^2(\mathbb R^d;\mathbb C)\) 上与可积 Fourier 变换一致。
\end{theorem}
\begin{proof}
对 \(f\in L^2\)，取 \(f_n=f\cdot\mathbf1_{\{|x|\leq n,\ |f|\leq n\}}\)。它有界、有限测度支集，故属于 \(L^1\cap L^2\)，并由控制收敛在 \(L^2\) 中趋于 \(f\)。Plancherel 等式说明 \((\widehat f_n)\) 是 Cauchy 列，定义 \(\mathcal F_2f\) 为其 \(L^2\) 极限。

若使用另一列 \(g_n\in L^1\cap L^2\) 在 \(L^2\) 中逼近 \(f\)，则
\[
 \|\widehat f_n-\widehat g_n\|_2
 =\|f_n-g_n\|_2
 \leq\|f_n-f\|_2+\|g_n-f\|_2\longrightarrow0.
\]
所以定义不依赖选择。特别当 \(f\) 本来属于 \(L^1\cap L^2\) 时，可以选恒定列 \(g_n=f\)，立刻得到新定义与原来的可积变换一致。

对 \(f,h\in L^2\) 分别取逼近列 \(f_n,h_n\)。列 \(af_n+bh_n\) 在交空间内并逼近 \(af+bh\)，因此积分变换的线性性经极限给出延拓的线性性。又有
\[
 \|\mathcal F_2f\|_2
 =\lim_n\|\widehat f_n\|_2
 =\lim_n\|f_n\|_2=\|f\|_2.
\]
任何连续延拓在稠密子集上相等，就在整个空间相等，故唯一。整个构造只使用 \(L^2\) 的完备性，并没有对不同频率逐点选取极限。
\end{proof}
\symidx{\(\mathcal F_2\)：平方可积函数的 Fourier 变换}

当不致混淆时仍记 \(\widehat f=\mathcal F_2f\)。例如一维函数 \((1+|x|)^{-3/4}\) 平方可积却不可积；对它不能用原来的 Lebesgue 积分定义变换，因为振荡因子模为 \(1\)，不改变绝对可积性。它的变换由上述范数极限定义。

\ProseParagraphBreak

若某个具体逼近同时在 \(L^1\) 和 \(L^2\) 中收敛，则两种极限也可直接核对。可积变换给出 \(\|\widehat f_n-\widehat f\|_\infty\leq\|f_n-f\|_1\)；延拓给出 \(\widehat f_n\to\mathcal F_2f\) 于 \(L^2\)。在任意有限球上，一致收敛意味着依测度收敛，而平方范数收敛也意味着依测度收敛，极限唯一。因此二者几乎处处相等。后面处理不同指数的交空间时，仍将使用这个识别步骤。

\subsection{酉性}
\label{b2:subsec:040103}

\begin{theorem}[Plancherel，酉形式]\label{b2:thm:04-unitary}
设 \(d\) 为正整数，\(\mathcal F_2:L^2(\mathbb R^d;\mathbb C)\to L^2(\mathbb R^d;\mathbb C)\) 是积分 \(\widehat f(\xi)=\int f(x)e^{-2\pi\mathrm ix\cdot\xi}\,\mathrm dx\) 从 \(L^1(\mathbb R^d;\mathbb C)\cap L^2(\mathbb R^d;\mathbb C)\) 的等距延拓。映射 \(\mathcal F_2\) 满射并保持内积，因而为\term[you-suanzi]{酉算子}[unitary operator]。若 \(Rf(x)=f(-x)\)，则
\[
 \mathcal F_2^2=R,\quad \mathcal F_2^4=I,
 \quad \mathcal F_2^{-1}=R\mathcal F_2=\mathcal F_2R.
\]
\end{theorem}
\begin{proof}
若 \(f\in C_c^\infty(\mathbb R^d;\mathbb C)\)，取整数 \(m>d\)，\cref{b2:prop:03-higher-derivative-decay}给出 \(\widehat f\) 在 \(|\xi|\geq1\) 上的 \(|\xi|^{-m}\) 衰减；在单位球内它又有界，所以 \(\widehat f\in L^1\)。它属于 \(L^2\) 则由已证的交空间 Plancherel 定理保证。因此可用反演得到 \(\mathcal F_2^2f=Rf\)。光滑紧支集函数在 \(L^2\) 中稠密，且两边都是连续算子，故等式推广到全部 \(L^2\)。于是 \(\mathcal F_2^4=R^2=I\)，特别每个 \(g\) 都是 \(\mathcal F_2(\mathcal F_2^3g)\)，所以满射。线性等距映射保持内积：将 \(f+g\)、\(f-g\)、\(f+\mathrm{i}g\)、\(f-\mathrm{i}g\) 的范数平方展开并作差，可分别恢复内积的实部和虚部。最后由平方等式及反射的交换关系得到逆算子公式；交换关系先由可积函数换元验证，再按稠密性延拓。
\end{proof}

在本书的归一化下，Fourier 变换没有丢失或放大总能量，但一般不保持某个点的值，也不保持函数的支集。酉性属于 Hilbert 空间结构，而不是逐点大小不变。

\subsection{\texorpdfstring{\(L^2\)}{L²} 反演与频率截断}
\label{b2:subsec:040104}

\begin{corollary}\label{b2:cor:04-l2-inversion}
设 \(d\) 为正整数，\(f\in L^2(\mathbb R^d;\mathbb C)\)，记 \(\widehat f=\mathcal F_2f\) 为其 Plancherel 变换。则
\[
 f=\lim_{R\to\infty}\int_{|\xi|\leq R}\widehat f(\xi)\cdot e^{2\pi\mathrm{i}x\cdot\xi}\,\mathrm{d}\xi
\]
在 \(L^2\) 范数中成立。当 \(\varepsilon\downarrow0\) 时，Gauss 正则化也有 \(\mathcal F_2^{-1}(e^{-\pi\varepsilon^2|\xi|^2}\widehat f)\to f\) 于 \(L^2\)。
\end{corollary}
\begin{proof}
记 \(h_R=\mathbf1_{\{|\xi|\leq R\}}\widehat f\)。Cauchy--Schwarz 给出
\[
 \|h_R\|_1\leq |B(0,R)|^{1/2}\|\widehat f\|_2,
\]
所以 \(h_R\in L^1\cap L^2\)。由延拓在交空间上的一致性和逆算子公式，\(\mathcal F_2^{-1}h_R\) 正是所写的正号积分；它有连续代表，但不要求 \(f\) 也预先连续。酉性给出
\[
 \|\mathcal F_2^{-1}h_R-f\|_2^2
 =\|h_R-\widehat f\|_2^2
 =\int_{|\xi|>R}|\widehat f(\xi)|^2\,\mathrm{d}\xi\longrightarrow0.
\]
Gauss 情形的乘子模不超过 \(1\)，并逐点趋于 \(1\)，故频率侧的控制收敛和酉性给出同一结论。
\end{proof}

任何递增穷尽空间的可测频率区域都可替代球；所需条件只是被舍弃部分的平方质量趋零。这个结论没有保证每个点上的截断积分收敛。对于随时间演化的方程，只要频率乘子有界，就可以先在 \(L^2\) 中定义演化，再另行研究点值和正则性。



\subsection{Sobolev 空间的频率刻画}
\label{b2:subsec:040105}

反演恢复的是平方范数意义下的函数。接下来进一步问：怎样从频率一侧判断它是否具有平方可积的弱导数？上一章对可积函数证明了“空间求导变成频率乘法”；这里须把这一对应建立在整个 \(L^2\) 空间上。下文用 \(\widehat f\) 表示由平方范数延拓得到的 Plancherel 变换。

\begin{theorem}[平方可积弱导数的频率刻画]\label{b2:thm:04-weak-derivative}
设 \(d\) 为正整数，\(j\in\{1,\ldots,d\}\)，\(f\in L^2(\mathbb R^d;\mathbb C)\)，并记 \(\widehat f=\mathcal F_2f\)。它具有平方可积的第 \(j\) 个弱导数，当且仅当 \(\xi_j\widehat f(\xi)\in L^2(\mathbb R^d)\)。此时
\[
 \mathcal F_2(\partial_jf)=2\pi\mathrm{i}\xi_j\widehat f,
 \quad\|\partial_jf\|_2=2\pi\|\xi_j\widehat f\|_2.
\]
\end{theorem}
\begin{proof}
先记下两种双线性积分等式，所有变换在此均取 \(L^2\) 意义：
\[
 \int\widehat u\,v=\int u\,\widehat v,
 \quad \int u\,v=\int\widehat u\,\check v,
 \quad \check v=\mathcal F_2^{-1}v.
\]
第一式先对 \(u,v\in L^1\cap L^2\) 由 Fubini 证明，交换重积分的上界为 \(\|u\|_1\|v\|_1\)。用交空间同时逼近两者，Cauchy--Schwarz 和等距性允许两侧取极限。将第一式中的 \(v\) 替换为 \(\check v\)，便得到第二式。这里没有复共轭，与弱导数的双线性测试约定一致；它不同于带有复共轭的内积等式。

若 \(g=\partial_jf\in L^2\)，弱导数测试式原先只用于 \(C_c^\infty\)。对 \(\psi\in C_c^\infty\)，矩公式给出 \(\widehat\psi\in C^\infty\)；又因 \(\psi\) 和 \(\xi_j\psi\) 都在 \(L^1\cap L^2\) 中，交空间 Plancherel 定理使 \(\widehat\psi\)、\(\partial_j\widehat\psi=\widehat{-2\pi\mathrm i\xi_j\psi}\) 都属于 \(L^2\)。取固定光滑截断 \(\chi\)，在单位球上等于 \(1\)、在半径 \(2\) 的球外为零，令 \(\chi_R(x)=\chi(x/R)\)。将 \(\chi_R\widehat\psi\) 代入测试式，得到
\[
 \int g\chi_R\widehat\psi
 =-\int f\chi_R\partial_j\widehat\psi
   -\int f(\partial_j\chi_R)\widehat\psi.
\]
前两项分别由 \(g,f\in L^2\) 和测试函数的平方可积性取极限。最后一项的绝对值至多为
\[
 R^{-1}\|\partial_j\chi\|_\infty\|f\|_2\|\widehat\psi\|_2,
\]
故趋于零。因此也可将 \(\widehat\psi\) 用作测试函数。由于直接对紧支集积分求导有 \(\partial_j\widehat\psi=\widehat{-2\pi\mathrm{i}\xi_j\psi}\)，双线性等式给出
\[
 \int\widehat g\,\psi
 =\int g\,\widehat\psi
 =-\int f\,\partial_j\widehat\psi
 =\int\widehat f\,(2\pi\mathrm{i}\xi_j\psi).
\]
两边的局部可积函数对全部紧支集光滑函数配对相同，故几乎处处相等。这个局部唯一性可由磨光核在 Lebesgue 点的逼近直接得出。于是必要条件成立。

反之，设 \(g=\mathcal F_2^{-1}(2\pi\mathrm{i}\xi_j\widehat f)\)。对 \(\varphi\in C_c^\infty\)，用双线性 Plancherel 等式与正号变换的求导公式，
\[
 \int g\varphi
 =\int2\pi\mathrm{i}\xi_j\widehat f\,\check\varphi
 =-\int\widehat f\,\check{\partial_j\varphi}
 =-\int f\partial_j\varphi.
\]
所以 \(g\) 是弱导数。范数式由酉性得到。
\end{proof}

设 \(m\) 为非负整数。整数阶 Sobolev 空间 \(H^m(\mathbb R^d;\mathbb C)=W^{m,2}(\mathbb R^d;\mathbb C)\) 要求函数及所有 \(|\alpha|\leq m\) 的弱导数都属于 \(L^2\)。对这些导数逐阶使用上一定理，便得到频率乘子 \((2\pi\mathrm i\xi)^\alpha\)。若 \(f\in H^m(\mathbb R^d;\mathbb C)\)，按第一卷定义 21.1.2 中每个多重指标计一次的范数，有
\[
 \|f\|_{H^m}^2
 =\int\left(\sum_{|\alpha|\leq m}(2\pi)^{2|\alpha|}\xi^{2\alpha}\right)|\widehat f(\xi)|^2\,\mathrm{d}\xi.
\]
括号中的多项式与 \((1+4\pi^2|\xi|^2)^m\) 可用只依赖 \(d,m\) 的正数互相控制，但高阶时一般不数值相等：后者的多项式展开带有组合系数。因此用后一权重定义的范数与既有范数等价，却不能静默地把等价写成相等。

\ProseParagraphBreak

这个等价可以直接从多项式系数看出。令 \(t_j=4\pi^2\xi_j^2\)，则第一种权重为 \(P_m(t)=\sum_{|\alpha|\leq m}t^\alpha\)。展开 \((1+t_1+\cdots+t_d)^m\) 后，每个 \(t^\alpha\) 的系数都是一个正整数，且只有有限多个。因此
\[
 P_m(t)\leq(1+t_1+\cdots+t_d)^m\leq C_{d,m}P_m(t)
 \quad(t_j\geq0).
\]
例如 \(d=m=2\) 时，\(P_2=1+t_1+t_2+t_1^2+t_1t_2+t_2^2\)，而另一权重的三个一次或混合项分别带有系数 \(2\)。这正对应了多重指标范数与展开 Laplace 算子的系数差别。

反过来，若 \(f\in L^2(\mathbb R^d;\mathbb C)\) 且 \((1+4\pi^2|\xi|^2)^{m/2}\widehat f\in L^2\)，上述多项式比较使每个 \((2\pi\mathrm i\xi)^\alpha\widehat f\)、\(|\alpha|\leq m\)，都属于 \(L^2\)。当 \(m\geq1\) 时，先用一阶定理构造一阶弱导数，再对已构造的导数重复；每一步需要的频率乘子都已在上述 \(L^2\) 控制之中，故逐阶得到全部高阶弱导数。\(m=0\) 时则只是原来的 \(L^2\) 条件。因此这个加权可积条件不仅是必要的，也是 \(f\in H^m\) 的充分条件。

\subsection{应用：常系数椭圆方程与差商}
\label{b2:subsec:040106}

从这里开始，利用已建立的 Fourier 变换与 Sobolev 空间的对应，通过频率乘子构造常系数椭圆方程的解。先看一个严格正的频率符号：除以它的形式计算，需要同时说明所得函数属于哪个空间、方程如何成立以及解为何唯一。

\begin{proposition}[常系数预解方程]\label{b2:prop:04-resolvent-model}
设 \(d\) 为正整数，\(a>0\)、\(f\in L^2(\mathbb R^d;\mathbb C)\)。方程 \((a-\Delta)u=f\) 有唯一的 \(H^2(\mathbb R^d;\mathbb C)\) 弱解，满足
\[
 \widehat u(\xi)=\frac{\widehat f(\xi)}{a+4\pi^2|\xi|^2},
 \quad \|u\|_2\leq a^{-1}\|f\|_2,
 \quad \|\Delta u\|_2\leq\|f\|_2.
\]
另有 \(\|\nabla u\|_2\leq(2\sqrt a)^{-1}\|f\|_2\)，并有仅依赖 \(a,d\) 的 \(H^2\) 估计。
这里弱解按分布意义理解，即对每个 \(\varphi\in C_c^\infty(\mathbb R^d)\)，均有 \(\int u(a\varphi-\Delta\varphi)=\int f\varphi\)；对 \(H^2\) 解，这也等价于方程在 \(L^2\) 中成立。
\end{proposition}
\begin{proof}
分母至少为 \(a\)，所以右侧属于 \(L^2\)，可由酉逆变换定义 \(u\)。置 \(s=2\pi|\xi|\)。对于零、一、二阶乘子，分别有
\[
 \frac1{a+s^2}\leq\frac1a,
 \quad\frac{s}{a+s^2}\leq\frac1{2\sqrt a},
 \quad\frac{s^2}{a+s^2}\leq1.
\]
中间不等式来自 \(a+s^2\geq2\sqrt a\,s\)。每个 \((2\pi\xi)^\alpha/(a+s^2)\)、\(|\alpha|\leq2\)，都受这三个界之一控制。迭代弱导数刻画得到 \(u\in H^2\)，而平方积分给出所述估计；对梯度先把各分量的平方相加，恰得到分子 \(s^2\)。由 \(\widehat{\Delta u}=-4\pi^2|\xi|^2\widehat u\)，有
\[
 \mathcal F_2((a-\Delta)u)=\widehat f.
\]
酉变换的单射性使等式在 \(L^2\) 中成立，因而也在弱意义下成立。若另一个 \(H^2\) 解存在，对两解之差 \(v\) 取变换，得到 \((a+4\pi^2|\xi|^2)\widehat v=0\)。正分母使 \(\widehat v=0\)，再由单射性得到 \(v=0\)。
\end{proof}

假设 \(a>0\) 使分母在低频也有统一正下界，从而保证对每个 \(L^2\) 右端都能构造解。令 \(a\downarrow0\) 时，\(L^2\) 解算子范数界 \(a^{-1}\) 发散；仅凭 \(f\in L^2\) 不能把分母换成 \(4\pi^2|\xi|^2\) 并断言仍得到平方可积解。这不表示 \(a=0\) 时出现了非零的 \(L^2\) 齐次解：若 \(v\in H^2\) 满足 \(-\Delta v=0\)，则 \(\widehat v\) 在 \(\xi\ne0\) 处为零，而单点 \(\{0\}\) 的测度为零，故 \(v=0\)。这里失去的是对任意右端的有界求解保证。相比之下，二阶导数的界在这个极限中保持一致。区分解本身与其导数的估计，能看出低频与高频承担的不同作用。\chapref{b2:ch:05}将进一步把 \(a=1\) 时的分母解释成基本解的变换。

\ProseParagraphBreak

弱导数的频率刻画还把局部差商与高频尾部连接起来。差商可以直接从函数值计算，因此适合判断一个具体的平方可积函数是否具有有限的梯度能量。

\begin{proposition}[平方范数差商]\label{b2:prop:04-difference-quotients}
设 \(d\) 为正整数，\(f\in L^2(\mathbb R^d;\mathbb C)\)，\(j\in\{1,\ldots,d\}\)，\(e_j\) 为第 \(j\) 个坐标单位向量。对 \(h\in\mathbb R\setminus\{0\}\)，记
\[
 D_{j,h}f(x)\coloneqq\frac{f(x+he_j)-f(x)}h.
\]
若弱导数 \(\partial_jf\in L^2(\mathbb R^d)\)，则 \(\|D_{j,h}f\|_2\leq\|\partial_jf\|_2\)，并且当 \(h\to0\) 时，\(D_{j,h}f\to\partial_jf\) 于 \(L^2\)。反之，若存在 \(h_n\in\mathbb R\setminus\{0\}\)、\(h_n\to0\) 及 \(C\geq0\)，使 \(\sup_n\|D_{j,h_n}f\|_2\leq C\)，则 \(f\) 具有平方可积的第 \(j\) 个弱导数，且 \(\|\partial_jf\|_2\leq C\)。
\end{proposition}
\begin{proof}
平移公式先在 \(L^1\cap L^2\) 上由换元成立，再由平移的等距性和稠密性推广到 \(L^2\)。因此
\[
 \widehat {D_{j,h}f}(\xi)
 =\frac{e^{2\pi\mathrm{i}h\xi_j}-1}{h}\widehat f(\xi).
\]
由 \(e^{\mathrm{i}t}-1=\int_0^t\mathrm{i}e^{\mathrm{i}s}\,\mathrm{d}s\)，乘子的模不超过 \(2\pi|\xi_j|\)，且逐点趋于 \(2\pi\mathrm{i}\xi_j\)。若这个极限乘以 \(\widehat f\) 属于 \(L^2\)，酉性给出差商的范数界。差商与极限之差的平方由 \(4(2\pi\xi_j)^2|\widehat f|^2\) 控制，控制收敛使它在 \(L^2\) 中趋于零，再用弱导数定理识别极限。

反之，若给定序列的差商一致有界，Fatou 引理直接给出
\[
 \int(2\pi\xi_j)^2|\widehat f(\xi)|^2\,\mathrm{d}\xi
 \leq\liminf_n\int\left|\frac{e^{2\pi\mathrm{i}h_n\xi_j}-1}{h_n}\right|^2
       |\widehat f(\xi)|^2\,\mathrm{d}\xi\leq C^2.
\]
所以频率刻画的充分条件成立，给出弱导数和所述界。无需先从差商中提取弱收敛子列。
\end{proof}

在一维，帐篷函数 \(\tau(x)=(1-|x|)_+\) 虽在 \(0\) 处不可微，却属于 \(H^1\)。在 \((-1,0)\)、\((0,1)\) 分别分部积分，端点处由函数连续使边界项相消，得到弱导数在这两段分别为 \(1\)、\(-1\)，其余为零，平方范数为 \(2\)。它也等于 \(\mathbf1_{[-1/2,1/2]}*\mathbf1_{[-1/2,1/2]}\)，所以
\[
 \widehat\tau(\xi)=\left(\frac{\sin(\pi\xi)}{\pi\xi}\right)^2.
\]
于是 \(\xi\widehat\tau\in L^2\)，与弱导数计算一致。相邻斜率的跳跃允许一阶能量有限，但若把帐篷换成矩形 \(v=\mathbf1_{[-1/2,1/2]}\)，便有 \(\xi\widehat v=\sin(\pi\xi)/\pi\notin L^2\)。矩形的函数值跳跃阻止了平方可积的一阶弱导数。这个检验区分了函数的折角与函数本身的不连续。

\ProseParagraphBreak

对于正整数 \(d,m\)、\(f\in H^m(\mathbb R^d;\mathbb C)\) 和 \(R>0\)，高频截断还有定量误差。令 \(S_Rf\coloneqq\mathcal F_2^{-1}(\mathbf1_{\{|\xi|\leq R\}}\widehat f)\)。由 \(|\xi|\geq R\) 在尾部成立，
\[
 \|f-S_Rf\|_2^2
 \leq R^{-2m}\int|\xi|^{2m}|\widehat f(\xi)|^2\,\mathrm{d}\xi
 \leq C_{d,m}R^{-2m}\sum_{|\alpha|=m}\|\partial^\alpha f\|_2^2.
\]
最后一步是多项式 \((\xi_1^2+\cdots+\xi_d^2)^m\) 的展开与弱导数公式。因而每增加一阶平方可积导数，舍弃高频后的范数误差就得到额外一阶尺度衰减。这个估计不把弱导数转换成逐点光滑性；它描述的是总能量在高频中的分布。

% !TeX root = ../../Book2.tex
\section{Hausdorff--Young 不等式}
\label{b2:sec:0402}

变换已经有两个端点估计：可积函数变为有界函数，平方可积函数保持平方范数。中间指数的控制将通过一个连续参数把这两端连起来。调节函数幅值的幂，可以使这个参数在两端分别对应已知的范数；\secref{b2:sec:1704}还将把同一方法用于一般线性算子。

\subsection{已知端点与相容定义}
\label{b2:subsec:04-fourier-endpoints}
% 以下两个旧标签兼容曾分别介绍端点的旧稿；合并后均指本小节。
% 标签后缀不决定打印编号；新引用用语义标签，不把旧040202移给插值小节。
\label{b2:subsec:040201}
\label{b2:subsec:040202}

先把待连接的两个端点和它们的共同定义域列清楚。在本卷 \(e^{-2\pi\mathrm ix\cdot\xi}\) 的约定下，对 \(f\in L^1(\mathbb R^d)\)，积分三角不等式给出 \(\|\widehat f\|_\infty\leq\|f\|_1\)；对 \(f\in L^2(\mathbb R^d)\)，Plancherel 延拓给出 \(\|\mathcal F_2f\|_2=\|f\|_2\)。第一个端点的常数为 \(1\)，非负可积函数在零频率处取等；第二个端点则对每个平方可积函数保持范数。

\ProseParagraphBreak

这里的 \(L^1\) 端点只控制每个频率处的大小，不能给出整个 \(L^1\) 单位球的统一衰减率，见\cref{b2:prop:03-no-uniform-decay}。其局部解释是：把一个非负单位积分的 \(\varphi\in C_c^\infty(\mathbb R^d)\) 调制成 \(f_\eta(x)=e^{2\pi\mathrm ix\cdot\eta}\varphi(x)\)，并不改变 \(L^1\) 范数，却使 \(\widehat f_\eta(\eta)=1\)，而 \(|\eta|\) 可以任意大。这与\cref{b2:thm:03-riemann-lebesgue}对每个固定函数的趋零结论相容：这里允许输入随所测的频率改变。

\ProseParagraphBreak

连接两个端点还需要相容性：\cref{b2:thm:04-plancherel-extension}说明，对每个 \(f\in L^1\cap L^2\)，平方范数变换 \(\mathcal F_2f\) 与绝对积分定义的 \(\widehat f\) 几乎处处相同。因此下面是在改变同一个算子的输入指数。有限测度支集的简单函数同时属于这两个空间；对它们使用复参数改变各个非零取值的幅度，两个边界分别落入 \(L^1\) 与 \(L^2\)，便能在共同定义域内完成插值。所得估计之后再用于一般 \(L^p\) 函数的延拓。

\subsection{Riesz--Thorin 思想}
\label{b2:subsec:04-complex-interpolation}
% 保留旧标识的数学含义，不按当前小节显示序号重新分配。
\label{b2:subsec:040203}

先用三线引理控制复参数的变化。下面出现的整函数都是有限指数和，可以由处处收敛的幂级数直接理解。

\begin{lemma}[三线引理的有界形式]\label{b2:lem:04-three-lines}
设整函数 \(F:\mathbb C\to\mathbb C\) 在闭带 \(0\leq\operatorname{Re}z\leq1\) 上有界，并对每个 \(t\in\mathbb R\) 满足 \(|F(\mathrm{i}t)|\leq1\)、\(|F(1+\mathrm{i}t)|\leq1\)。则 \(|F(\theta)|\leq1\) 对 \(0\leq\theta\leq1\) 成立。
\end{lemma}
\begin{proof}
先回顾幂级数的最大模原理。若中心的最大模为零，函数已在邻域内为零；否则把函数乘以常数相位，使中心值为正实数 \(A\)。在足够小的圆上，幂级数一致绝对收敛，逐项平均合法，所得平均为中心值 \(A\)。圆上实部至多为模，模又至多为 \(A\)，所以连续非负函数 \(A-\operatorname{Re}F\) 的圆周积分为零，它处处为零。此时 \(|F|\leq A\) 还迫使虚部为零，故圆上函数恒为 \(A\)。将圆周值乘以各个负频率后逐项积分，可见所有正阶幂级数系数为零。函数在邻域内常值，并由重叠圆盘延拓到连通域。因此非恒定解析函数在紧矩形上的最大模出现在边界。

给定 \(\varepsilon>0\)，考虑 \(F(z)e^{\varepsilon z^2}\)。在两条竖边上，其模至多为 \(e^\varepsilon\)。若 \(|F|\leq M\)，在高度为 \(\pm R\) 的两条横边上，模至多为 \(Me^{\varepsilon(1-R^2)}\)，对充分大的 \(R\) 也不超过 \(e^\varepsilon\)。在矩形上用刚证明的最大模原理，得 \(|F(\theta)|e^{\varepsilon\theta^2}\leq e^\varepsilon\)。最后令 \(\varepsilon\downarrow0\)。
\end{proof}

这一估计是\term[sanxian-yinli]{三线引理}[three-lines lemma]的特殊形式。它把两条边界上的控制传递到中间，下面让边界对应两种不同的积分指数。

\begin{theorem}[Hausdorff--Young]\label{b2:thm:04-hausdorff-young}
设 \(d\) 为正整数，\(1\leq p\leq2\)，\(q=p'\)，即 \(1/p+1/q=1\)，其中 \(p=1\) 时 \(q=\infty\)。若 \(f\in L^1(\mathbb R^d;\mathbb C)\cap L^p(\mathbb R^d;\mathbb C)\)，则
\[
 \|\widehat f\|_q\leq\|f\|_p.
\]
\end{theorem}
\begin{proof}
端点已证，以下设 \(1<p<2\)，于是 \(2<q<\infty\) 且 \(q'=p\)。先设 \(f\) 为有界有限测度支集的简单函数。零函数直接满足结论；否则除以其 \(L^p\) 范数，先证明 \(\|f\|_p=1\) 的情形。再取任意同类简单函数 \(g\)，满足 \(\|g\|_p=1\)。我们尚不知道 \(\widehat f\in L^q\)，所以先估计配对 \(\int\widehat f g\)，随后在有限球上从这些配对恢复局部 \(L^q\) 范数，最后才扩大到全空间。测试函数的指数是 \(p=q'\)，因此可以对 \(f\) 和 \(g\) 使用相同的变幂构造。

现在说明幂指数如何由目标决定。设幅度的指数取实系数仿射函数 \(a(z)=A+Bz\)，因而在非零取值处，复幂 \(|f|^{a(z)}\) 的模为 \(|f|^{a(\operatorname{Re}z)}\)。要使两条边界恰好使用 \(L^1\) 与 \(L^2\) 端点，并在某个实数 \(0<\theta<1\) 处恢复原函数，必须同时满足
\[
 \begin{aligned}
  \operatorname{Re}z=0:&\quad a(0)=p,
       &&\text{使 }|f_z|=|f|^p,\quad\|f_z\|_1=1;\\
  \operatorname{Re}z=1:&\quad 2a(1)=p,
       &&\text{使 }|f_z|^2=|f|^p,\quad\|f_z\|_2=1;\\
  z=\theta:&\quad a(\theta)=1,
       &&\text{使 }f_\theta=f.
 \end{aligned}
\]
前两式确定 \(A=p\)、\(B=-p/2\)，第三式再确定 \(\theta=2(1-1/p)\)，在 \(1<p<2\) 时确有 \(0<\theta<1\)。这三个条件共同确定了指数 \(a(z)=p(1-z/2)\)；\(g\) 的指数也由相同条件确定。

记非零复数 \(w\) 的相位为 \(\operatorname{phase}(w)=w/|w|\)，零处定义为零。令
\[
 f_z(x)=\operatorname{phase}(f(x))\,|f(x)|^{p(1-z/2)},
 \quad g_z(\xi)=\operatorname{phase}(g(\xi))\,|g(\xi)|^{p(1-z/2)},
\]
零值处仍取零，复幂使用正数 \(|f|\) 和 \(|g|\) 的实对数。对每个 \(z\)，这两个函数仍有界并分别支于原来的有限测度集合，因而下式的绝对值二重积分为 \(\|f_z\|_1\cdot\|g_z\|_1<\infty\)，可按 Fubini 写为 \(\int\widehat {f_z}g_z\)。定义
\[
 F(z)=\int\!\int e^{-2\pi\mathrm{i}x\cdot\xi}f_z(x)\cdot g_z(\xi)\,\mathrm{d}x\,\mathrm{d}\xi
\]
因为函数各只有有限个非零取值，\(F\) 是有限指数和。具体地，将非零取值写成 \(f=\sum_j a_j\mathbf1_{E_j}\)、\(g=\sum_\ell b_\ell\mathbf1_{H_\ell}\)，其中各集合在相应变量中互不相交。令
\[
 c_{j\ell}=\int_{E_j}\int_{H_\ell}e^{-2\pi\mathrm{i}x\cdot\xi}\,\mathrm{d}\xi\,\mathrm{d}x,
 \quad |c_{j\ell}|\leq |E_j|\cdot|H_\ell|.
\]
则 \(F(z)\) 是 \(c_{j\ell}\operatorname{phase}(a_j)\operatorname{phase}(b_\ell)\exp(p(1-z/2)(\log|a_j|+\log|b_\ell|))\) 的有限和。因对数均为实数，每一项沿带的虚方向只改变相位，该项的模仅取决于 \(\operatorname{Re}z\)；对有限项使用三角不等式，便给出 \(F\) 在整个闭带上的统一上界。这同时验证了整性与三线引理要求的全带有界性；仅仅知道两条边上的估计，还不能省略后一条件。

\cref{b2:fig:04-interpolation-strip}把三个指数条件放在同一条带上：虚方向只改变相位，左右边界的幅度分别恢复归一化的 \(L^1\) 与 \(L^2\) 范数。下面逐一验证边界上的配对估计，再用三线引理传到 \(z=\theta\)。

\begin{figure}[htbp]
\centering
\begin{tikzpicture}[x=1cm,y=1cm,font=\small]
 \fill[AnalysisBlue!7] (0,-1.1) rectangle (9,1.1);
 \draw[->,AnalysisGray] (-0.5,0)--(9.7,0) node[right] {\(\operatorname{Re}z\)};
 \draw[<->,thick,AnalysisBlue] (0,-1.2)--(0,1.2);
 \draw[<->,thick,AnalysisBlue] (9,-1.2)--(9,1.2);
 \draw[<->,dashed,AnalysisBlue] (4.5,-1.2)--(4.5,1.2);
 \node[above left] at (0,1.2) {\(\operatorname{Im}z\)};
 \node[below left] at (0,0) {\(0\)};
 \fill[AnalysisBlue] (4.5,0) circle (1.5pt);
 \node[below right] at (4.5,0) {\(\theta\)};
 \node[below right] at (9,0) {\(1\)};
 \node[align=center,anchor=south] at (0,1.65)
   {\(\|f_{\mathrm it}\|_1=\|g_{\mathrm it}\|_1=1\)\\\(L^1\longrightarrow L^\infty\)};
 \node[align=center,anchor=south] at (4.5,1.65)
   {\(f_\theta=f,\quad g_\theta=g\)\\\(\|f\|_p=\|g\|_p=1\)};
 \node[align=center,anchor=south] at (9,1.65)
   {\(\|f_{1+\mathrm it}\|_2=\|g_{1+\mathrm it}\|_2=1\)\\\(L^2\longrightarrow L^2\)};
 \node[below] at (0,-1.35) {\(|F(\mathrm it)|\leq1\)};
 \node[below] at (9,-1.35) {\(|F(1+\mathrm it)|\leq1\)};
 \draw[->,AnalysisBlue] (0.8,-0.65)--(3.35,-0.65);
 \draw[->,AnalysisBlue] (8.2,-0.65)--(5.65,-0.65);
 \node[fill=white,inner sep=3pt] at (4.5,-0.65) {\(|F(\theta)|\leq1\)};
\end{tikzpicture}
\caption[复插值带与端点范数]{复插值带与端点范数。这里只画出无限高带状区域的一段，\(\theta\) 的横向位置仅作示意；两条边界上的估计由三线引理传递到中间。}
\label{b2:fig:04-interpolation-strip}
\end{figure}

当 \(z=\mathrm{i}t\) 时，\(|f_z|=|f|^p\)，故 \(\|f_z\|_1=\int|f|^p=1\)，\(g_z\) 同理。于是
\[
 |F(\mathrm{i}t)|\leq\|\widehat {f_{\mathrm{i}t}}\|_\infty\cdot\|g_{\mathrm{i}t}\|_1\leq1.
\]
当 \(z=1+\mathrm{i}t\) 时，\(|f_z|^2=|f|^p\)，所以 \(\|f_z\|_2=\|g_z\|_2=1\)。利用交空间上的 Plancherel 等式，
\[
 |F(1+\mathrm{i}t)|\leq\|\widehat {f_{1+\mathrm{i}t}}\|_2\cdot\|g_{1+\mathrm{i}t}\|_2=1.
\]
三线引理给出 \(|F(\theta)|\leq1\)；而 \(p(1-\theta/2)=1\)，所以
\[
 \left|\int\widehat f(\xi)g(\xi)\,\mathrm{d}\xi\right|\leq1.
\]
下面从配对估计恢复范数。若直接写全空间 \(L^q\) 的对偶表示，就预先假设了正在证明的 \(\widehat f\in L^q\)。改在有限球上测试可以避开这一点：记 \(h=\widehat f\)，由 \(f\in L^1\) 已知 \(h\) 有界，故对每个 \(R>0\)，它自动属于 \(L^q(B_R)\)，即使尚不知道全空间的积分是否有限。令
\[
 A_R=\left(\int_{B_R}|h|^q\right)^{1/q}.
\]
若 \(A_R=0\) 则已有 \(A_R\leq1\)。否则，为使配对恰好恢复 \(A_R\)，按 Hölder 取等的形式选测试函数：其相位取 \(h\) 相位的共轭以消去振荡，幅度取 \(|h|^{q-1}\)，再除以使 \(L^p\) 范数为 \(1\) 的常数。具体取
\[
 g_R=\frac{\overline{\operatorname{phase}(h)}\,|h|^{q-1}\mathbf1_{B_R}}{A_R^{q-1}}.
\]
因为 \((q-1)p=q\)，直接计算得 \(\|g_R\|_p=1\)、\(\int hg_R=A_R\)。函数 \(g_R\) 有界并支于有限球，可以用同一球上的简单函数在 \(L^p\) 中逼近。此时 \(h\in L^q(B_R)\)，Hölder 不等式保证配对收敛；逼近函数的范数也趋于 \(1\)，归一化后可用已经证明的简单函数估计。因此 \(A_R\leq1\)。最后令 \(R\uparrow\infty\)，单调收敛给出 \(\int|h|^q\leq1\)。去掉最初的归一化，便得简单函数的结论。

最后设 \(f\in L^1\cap L^p\)。先同时截断空间与幅值，再对截断函数量化取值，可得到简单函数 \(s_n\) 同时在 \(L^1\) 与 \(L^p\) 中趋于 \(f\)。第一个收敛使 \(\widehat s_n\) 一致趋于普通变换 \(\widehat f\)，第二个使 \(\|s_n\|_p\to\|f\|_p\)。Fatou 引理遂给出
\[
 \int|\widehat f|^q\leq\liminf_n\int|\widehat s_n|^q
 \leq\lim_n\|s_n\|_p^q=\|f\|_p^q.
\]
因此不等式对每个 \(f\in L^1\cap L^p\) 成立。
\end{proof}

上述证明体现\term[riesz-thorin-chazhi]{Riesz--Thorin 插值}[Riesz--Thorin interpolation]的方法：变动的是函数幅值的幂，使两个边界恰好落在已知空间。结果称为\term[hausdorff-young-budengshi]{Hausdorff--Young 不等式}[Hausdorff--Young inequality]。这里的估计取常数 \(1\)；最佳常数问题还要比较所有非零输入的范数比值。

\subsection{\texorpdfstring{\(1\leq p\leq2\)}{1≤p≤2}}
\label{b2:subsec:04-lp-fourier-extension}
\label{b2:subsec:040204}

\begin{corollary}\label{b2:cor:04-lp-extension}
设 \(d\) 为正整数，\(1<p<2\)，\(p'=p/(p-1)\)，各函数空间均为 \(\mathbb R^d\) 上的复值 Lebesgue 空间。Fourier 变换唯一延拓为连续线性映射 \(\mathcal F_p:L^p(\mathbb R^d;\mathbb C)\to L^{p'}(\mathbb R^d;\mathbb C)\)，且 \(\|\mathcal F_pf\|_{p'}\leq\|f\|_p\)。它在 \(L^1\cap L^p\) 上与可积变换一致，在 \(L^2\cap L^p\) 上与 Plancherel 变换一致。
\end{corollary}
\begin{proof}
取有界有限测度支集的简单函数 \(s_n\to f\) 于 \(L^p\)。前一定理作用于差，给出
\[
 \|\widehat s_n-\widehat s_m\|_{p'}\leq\|s_n-s_m\|_p.
\]
故变换是 \(L^{p'}\) Cauchy 列，以其极限定义 \(\mathcal F_pf\)。两个逼近列的差趋于零，使定义独立于选择；对线性组合选择相应逼近列，再取范数极限，得到线性性及范数界。稠密性保证连续延拓唯一。

若 \(f\in L^p\cap L^r\)、\(r=1\) 或 \(2\)，\cref{b2:prop:03-common-smooth-density}提供同时在两个范数中逼近的 \(C_c^\infty\) 列 \(v_n\)。它们的变换在 \(L^{p'}\) 中趋于 \(\mathcal F_pf\)，在 \(r=1\) 时一致趋于可积变换，在 \(r=2\) 时以平方范数趋于 \(\mathcal F_2f\)。对任意有限球 \(B\)，Chebyshev 不等式使每种有限指数范数收敛都推出依测度收敛，例如
\[
 |\{\xi\in B:|\widehat v_n-\mathcal F_pf|>\delta\}|
 \leq\delta^{-p'}\|\widehat v_n-\mathcal F_pf\|_{p'}^{p'}\longrightarrow0.
\]
若两个极限分别记为 \(H_1,H_2\)，则集合 \(\{|H_1-H_2|>2\delta\}\cap B\) 包含于相应两个误差集合的并，因而测度为零。令 \(\delta\downarrow0\)，再让有限球穷尽空间，得到几乎处处相等。这完成了交空间的一致性核对。
\end{proof}

因此在共同定义域上可以继续使用 \(\widehat f\) 这个记号。若 \(1<p<2\)，变换属于 \(L^{p'}\)，但 \(p'>2\)，不能简单地再对变换使用同一个 Hausdorff--Young 定理来反演；反演的正确表达要看它是否另外满足可积性、平方可积性，或采用后面的分布框架。



\ProseParagraphBreak

一个新定义的必要性，可以从同时不属于两个旧空间的函数看出。固定 \(1<p<2\)，选
\[
 d/2<\alpha<d/p<\beta<d,
\]
并取 \(f(x)=|x|^{-\alpha}\mathbf1_{\{|x|<1\}}+(1+|x|)^{-\beta}\mathbf1_{\{|x|\geq1\}}\)。原点附近由 \(\alpha p<d\) 保证 \(p\) 次可积，无穷远由 \(\beta p>d\) 保证，所以 \(f\in L^p\)。但原点奇性使它不平方可积，无穷远尾部使它不可积。因而它既不能整体代入\chapref{b2:ch:03}的绝对积分，也不能整体代入 Plancherel 定义；Hausdorff--Young 延拓确实增加了对象。

\ProseParagraphBreak

具体的积分检验是 \(\int_0^1 r^{d-1-\alpha p}\,\mathrm{d}r<\infty\) 与 \(\int_1^\infty r^{d-1-\beta p}\,\mathrm{d}r<\infty\)。把指数分别换成 \(2\) 与 \(1\)，第一个积分因 \(2\alpha>d\) 发散，第二个因 \(\beta<d\) 发散。例如一维取 \(p=3/2\)、\(\alpha=3/5\)、\(\beta=4/5\)，就得到可直接检验的具体函数。这里原点与远处是两种不同的障碍，仅靠一次空间截断无法同时移除；逼近时需要空间截断和幅值截断共同作用。

\ProseParagraphBreak

定义这个变换时，只需取在 \(L^p\) 中逼近 \(f\) 的简单函数列。其频率像由不等式组成 \(L^{p'}\) Cauchy 列，不必事先猜测极限的点值公式。若另外将 \(f\) 拆成可积部分与平方可积部分，两个变换的和也应与新定义相容；这种相容性将在下一章的测试配对中统一表达。

% !TeX root = ../../Book2.tex
\section{\texorpdfstring{\(L^p\)}{Lᵖ} Fourier 变换的局限}
\label{b2:sec:0403}

把指数从 \(2\) 往两侧移动，会改变变换所能取值的空间。下面用尺度和相互分离的波包解释为什么 \(p>2\) 不能照搬上一节的结论，并为分布意义下的变换留下明确的问题。

\subsection{\texorpdfstring{\(p>2\)}{p>2} 时的有界延拓障碍}
\label{b2:subsec:040301}

先作尺度检验。假设存在 \(\|\widehat f\|_q\leq C\|f\|_p\) 对全部光滑紧支集函数成立。固定非零 \(\varphi\in C_c^\infty\)，其变换快速下降，故所需的各个范数均有限且非零。对 \(f_a(x)=\varphi(ax)\)、\(a>0\)，伸缩公式给出
\[
 \|f_a\|_p=a^{-d/p}\|\varphi\|_p,
 \quad\|\widehat {f_a}\|_q=a^{-d+d/q}\|\widehat\varphi\|_q,
\]
无穷指数按倒数为零解释。将两式代入假定的估计并消去固定正数，得到
\[
 a^{d(1/p+1/q-1)}\leq C_\varphi\quad(a>0).
\]
若幂为正，让 \(a\to\infty\)；若幂为负，让 \(a\downarrow0\)，均产生矛盾。因此必须有 \(1/p+1/q=1\)。尺度检验只筛去不可能的指数，并不证明剩下的指数可用。

\begin{theorem}\label{b2:thm:04-p-above-two}
设 \(d\) 为正整数，\(2<p\leq\infty\)、\(1/p+1/q=1\)，其中 \(1/\infty=0\)。采用 \(\widehat f(\xi)=\int_{\mathbb R^d}f(x)e^{-2\pi\mathrm{i}x\cdot\xi}\,\mathrm{d}x\) 的归一化，则不存在与 \(f\) 无关的有限常数 \(C\)，使 \(\|\widehat f\|_q\leq C\|f\|_p\) 对所有复值函数 \(f\in C_c^\infty(\mathbb R^d;\mathbb C)\) 成立。
\end{theorem}
\begin{proof}
固定非零 \(\varphi\in C_c^\infty\)，选 \(L\) 足够大，使 \(\varphi(x-jLe_1)\) 的支集互不相交。对符号 \(\epsilon_j\in\{-1,1\}\)，令
\[
 f_N(x)=\sum_{j=1}^N\epsilon_j\varphi(x-jLe_1).
\]
于是 \(\|f_N\|_p=N^{1/p}\|\varphi\|_p\)，\(p=\infty\) 时取 \(N^{1/p}=1\)；其变换为 \(\widehat\varphi(\xi)\sum_{j=1}^N\epsilon_je^{-2\pi\mathrm{i}jL\xi_1}\)。

把全部 \(2^N\) 种符号等权平均。这里 \(\mathbb E\) 只是有限和除以 \(2^N\)，\(\mathbb P\) 只是所占符号组的比例。对任意模为 \(1\) 的复数 \(z_j\)，记 \(Z=\sum_j\epsilon_jz_j\)、\(X=|Z|^2\)。翻转某个符号会使含它奇数次的单项变号，因此这种单项的平均为零。于是
\[
 \mathbb E|Z|^2=\sum_{j,k}z_j\overline{z_k}\,\mathbb E(\epsilon_j\epsilon_k)=N.
\]
展开 \(|Z|^4\) 时，一个四重指标只有每个指标出现偶数次才能保留。它必属于三个成对划分之一：第一位置与第二、第一与第三、或第一与第四配对。每种划分至多有 \(N^2\) 个项，每项系数模为 \(1\)。全相等的指标可能被多次计入，但我们只求上界，故 \(\mathbb EX^2\leq3N^2\)。令 \(E=\{X\geq N/2\}\)，在其补集上 \(X<N/2\)，而在 \(E\) 上使用 Cauchy--Schwarz，得
\[
 N=\mathbb EX\leq N/2+(\mathbb EX^2)^{1/2}\mathbb P(E)^{1/2},
\]
因此 \(\mathbb P(E)\geq1/12\)，进而 \(\mathbb E|Z|^q\geq12^{-1}(N/2)^{q/2}\)。

逐频率取 \(z_j=e^{-2\pi\mathrm{i}jL\xi_1}\)，乘以 \(|\widehat\varphi(\xi)|^q\) 再积分。有限平均可与积分交换，得到
\[
 2^{-N}\sum_{\epsilon\in\{-1,1\}^N}\|\widehat {f_N}\|_q^q
 \geq12^{-1}(N/2)^{q/2}\|\widehat\varphi\|_q^q.
\]
至少一个被平均的非负数不小于平均值，故有一组选定符号使
\[
 \|\widehat {f_N}\|_q\geq c_qN^{1/2}\|\widehat\varphi\|_q.
\]
光滑紧支集函数的变换快速衰减且非零，故右侧范数有限而正。若所求估计成立，则 \(N^{1/2-1/p}\) 应一致有界，与 \(p>2\) 矛盾。
\end{proof}

这个反例甚至使用光滑紧支集函数，故障碍来自范数估计本身，并非某个积分未定义。它否定的是在整个 \(L^p\) 上取值于 \(L^{p'}\) 的有界延拓；并没有否定特殊函数仍可能有很好的变换。

\subsection{普通积分定义的局限}
\label{b2:subsec:040302}

在无限测度空间上，较大的 \(p\) 不意味着 \(L^1\)。例如在 \(\mathbb R^d\) 上，对 \(2<p<\infty\) 取 \(d/p<\alpha\leq d\)，则 \((1+|x|)^{-\alpha}\in L^p\setminus L^1\)；对 \(p=\infty\)，任取 \(0<\alpha\leq d\) 也得到有界但不可积的函数。乘上单位模复指数以后，绝对可积性没有改善，所以定义式不能作为普通 Lebesgue 积分使用。这并不等于某种指定的振荡截断必定发散，而是说普通积分尚未提供定义。

\ProseParagraphBreak

有时可以指定对称截断、Gauss 正则化或其他极限。不同极限的存在性各需证明；一般的 \(L^p\) 等价类也没有预先规定的点值。即使在 \(L^2\) 中，构造的是频率侧的等价类；一句“取每个频率的极限”并不能替代范数完备性给出的定义。

\subsection{\texorpdfstring{\(L^2\)}{L²} Fourier 乘子及其推广障碍}
\label{b2:subsec:040303}

\begin{definition}\label{b2:def:04-multiplier}
设 \(d\) 为正整数。对 \(m\in L^\infty(\mathbb R^d;\mathbb C)\)、\(f\in L^2(\mathbb R^d;\mathbb C)\)，以 \(2\pi\) 归一化的 Plancherel 变换 \(\widehat f=\mathcal F_2f\) 定义
\[
 T_mf\coloneqq\mathcal F_2^{-1}(m\widehat f),
\]
称为符号 \(m\) 的\term[fourier-chengzi-suanzi]{Fourier 乘子算子}[Fourier multiplier operator]。
\end{definition}
\symidx{\(T_m\)：符号为 \(m\) 的 Fourier 乘子算子}

\begin{proposition}\label{b2:prop:04-multiplier-norm}
设 \(d\) 为正整数，\(\mathcal F_2\) 为 \(2\pi\) 归一化的 Plancherel 变换，\(m\in L^\infty(\mathbb R^d;\mathbb C)\)，\(T_mf=\mathcal F_2^{-1}(m\mathcal F_2f)\) 定义于 \(L^2(\mathbb R^d;\mathbb C)\)。有 \(\|T_m\|_{L^2\to L^2}=\|m\|_\infty\)。若 \(K\in L^1(\mathbb R^d;\mathbb C)\)、\(m=\widehat K\) 几乎处处成立，则 \(T_mf=K*f\) 对每个 \(f\in L^2(\mathbb R^d;\mathbb C)\) 成立，等式按 \(L^2\) 等价类理解；卷积公式还定义了每个 \(L^p(\mathbb R^d;\mathbb C)\)、\(1\leq p\leq\infty\) 上的有界算子，范数至多为 \(\|K\|_1\)。
\end{proposition}
\begin{proof}
酉性给出上界。当 \(\|m\|_\infty=0\) 时两边均为零；否则，对任意 \(0<a<\|m\|_\infty\)，集合 \(\{|m|>a\}\) 有正测度。用有限球穷尽空间，可在其中取一个有限正测度子集 \(E\)。由满射性可选 \(f\in L^2\) 使 \(\widehat f=\mathbf1_E\)，于是 \(\|T_mf\|_2>a\|f\|_2\)，令 \(a\uparrow\|m\|_\infty\) 得到等号。卷积表示先对 \(f\in L^1\cap L^2\) 由\cref{b2:thm:03-convolution-transform}成立：Young 不等式保证 \(K*f\) 同时属于 \(L^1\) 与 \(L^2\)，从而可积变换与 Plancherel 变换的一致性可用于两侧。再取 \(f_n\in L^1\cap L^2\) 在 \(L^2\) 中逼近任意 \(f\)；Young 估计使 \(K*f_n\to K*f\)，已证乘子界使 \(T_mf_n\to T_mf\)，故等式推广到 \(L^2\)。其余指数直接由\cref{b2:thm:03-young}给出。
\end{proof}

对 \(p\ne2\)，仅有 \(m\) 有界通常不足以取得同样的乘子结论；符号的平滑程度、奇性与维数会参与估计。\chapref{b2:ch:22}将研究这些性质怎样决定乘子的有界指数范围。


\ProseParagraphBreak

平滑频率截断说明为何符号的导数有用。取 \(m\in C_c^\infty(\mathbb R^d;\mathbb C)\)，令 \(K=\check m\)。\cref{b2:prop:03-higher-derivative-decay}将可积的高阶纯坐标导数转成频率衰减；选取大于维数的阶数，可得 \(K\in L^1\)，所以 \(T_m\) 在每个 \(L^p\) 上都由 Young 不等式控制。对 \(R>0\) 令 \(m_R(\xi)=m(\xi/R)\)，相应核为 \(K_R(x)=R^dK(Rx)\)，且 \(\|K_R\|_1=\|K\|_1\)。于是这一族算子对全部尺度具有同一个范数界。

\begin{proposition}[平滑频率逼近]\label{b2:prop:04-smooth-frequency-approximation}
设 \(d\) 为正整数，\(m\in C_c^\infty(\mathbb R^d;\mathbb C)\)、\(m(0)=1\)，按正号 Fourier 积分定义 \(K=\check m\)，并对每个 \(R>0\) 令 \(K_R(x)=R^dK(Rx)\)。对 \(1\leq p<\infty\)、\(f\in L^p(\mathbb R^d;\mathbb C)\)，当 \(R\to\infty\) 时有 \(K_R*f\to f\) 于 \(L^p\)；若复值函数 \(f\) 在 \(\mathbb R^d\) 上有界且一致连续，则收敛为一致收敛。对 \(f\in L^2(\mathbb R^d;\mathbb C)\)，这个卷积等于 \(T_{m(\cdot/R)}f\)。
\end{proposition}
\begin{proof}
先核对核的可积性。取整数 \(N>d\)，则 \(m\) 及各 \(\partial_j^Nm\) 都属于 \(L^1\)；由\cref{b2:prop:03-higher-derivative-decay}以及 \(\check m(x)=\widehat m(-x)\)，有 \(|K(x)|\leq C|x|^{-N}\) 对 \(|x|\geq1\) 成立，而在单位球内 \(|K|\leq\|m\|_1\)。两部分合起来给出 \(K\in L^1\)。对光滑紧支集的 \(m\) 使用上一章的反演，再作反射换元，得到 \(\widehat K=m\)；这里两边连续，所以等式处处成立。特别 \(\int K=\widehat K(0)=m(0)=1\)。注意 \(K\) 可以有复值且不必非负，所需的是绝对值的统一质量界与尾部集中性。将卷积换元，得到
\[
 K_R*f-f=\int K(y)\bigl(f(\cdot-y/R)-f\bigr)\,\mathrm{d}y.
\]
因而由积分型 Minkowski 不等式，
\[
 \|K_R*f-f\|_p
 \leq\int|K(y)|\,\|f(\cdot-y/R)-f\|_p\,\mathrm{d}y.
\]
对每个固定 \(y\)，有限指数的平移连续性使内层范数趋于零；它又不超过 \(2\|f\|_p\)。控制收敛给出结论。若 \(f\) 有界一致连续，用同一论证把 \(L^p\) 范数换成上确界范数即可。\(L^2\) 上的乘子识别来自前一命题，因为 \(\widehat {K_R}(\xi)=m(\xi/R)\)。
\end{proof}

这个证明也显示为何不能把结论无条件推广到 \(L^\infty\) 范数。在一维取阶跃函数 \(f=\mathbf1_{(0,\infty)}\)，它与光滑核卷积后具有连续代表。若一个连续函数与 \(f\) 的本性上确界距离小于 \(1/2\)，可取统一的 \(c<1/2\)，使它在左侧几乎处处距 \(0\) 至多为 \(c\)，在右侧几乎处处距 \(1\) 至多为 \(c\)。连续性使这些闭圆盘约束在原点也同时成立，但两圆盘不相交。因此 \(\|K_R*f-f\|_\infty\geq1/2\)；这里的收敛限制由函数的跳跃造成，尽管每个算子的 \(L^\infty\) 范数都统一有界。

\ProseParagraphBreak

相反，对 \(R>0\)，一维的硬截断符号 \(m=\mathbf1_{[-R,R]}\) 对应
\[
 K(x)=\dfrac{\sin(2\pi Rx)}{\pi x},
\]
该核不属于 \(L^1\)。对 \(f\in C_c^\infty\)，仍可用有限频率积分和 Fubini 验证 \(T_mf=K*f\)，但不能再用 \(\|K\|_1\) 给出有限的 Young 上界。\(L^2\) 上的范数仍为 \(1\)，因为符号是有界的正交投影。这个例子把两种方法的职责分开：酉性直接控制平方能量，绝对可积核则提供全部指数上的一种充分条件。

\subsection{分布意义的 Fourier 变换的动机}
\label{b2:subsec:040304}

若不能要求变换是一个普通函数，还可以要求它对足够好的测试函数产生确定的数值。对 \(f\in L^p\) 和光滑紧支集的 \(\varphi\)，函数 \(\widehat\varphi\) 快速衰减，因而 \(\int f\widehat\varphi\) 总有意义。若 \(f\) 原来可积，Fubini 又给出
\[
 \int\widehat f(\xi)\varphi(\xi)\,\mathrm{d}\xi
 =\int f(x)\widehat\varphi(x)\,\mathrm{d}x.
\]
右侧因此提供了延拓变换的线索，且对每个 \(1\leq p\leq\infty\) 都适用。下一章将选取在 Fourier 变换下保持不变的 Schwartz 测试空间，使这个配对成为完整定义。

\ProseParagraphBreak

测试函数的下降速度在这里承担了可积性责任。例如取整数 \(N\) 满足 \(Np'>d\)，当 \(p'=\infty\) 时只需取任意正整数。则 Hölder 不等式给出
\[
 \left|\int f\widehat\varphi\right|
 \leq C_{N,p,d}\|f\|_p
 \sup_x(1+|x|)^N|\widehat\varphi(x)|.
\]
其中 \(p'=1\) 也由 \(N>d\) 包含在内。反复分部积分把右侧的加权上确界控制在 \(\varphi\) 的有限多个导数与紧支集大小之内。这样一个可能没有可积频率代表的对象，仍能对每个测试函数给出稳定的线性数值；下一章要精确说明的是这种稳定性应在哪个测试函数拓扑中表述。


\begin{chaptersummary}
Plancherel 等式把变换从绝对收敛积分延拓为能量空间上的酉算子。反演和四次变换等于恒等的关系保证满射性；频率截断的误差则由被舍弃的平方质量精确给出。Hausdorff--Young 把两个端点连接起来，但其共轭指数方向只覆盖 \(1\leq p\leq2\)。

\ProseParagraphBreak

大于二时，空间中相互分离的函数可在频率中产生更大的叠加，破坏所需范数界。应当区分某个好函数的变换存在、变换在整个空间上连续，以及逐点反演这三个层次。\(L^2\) 乘子由符号本性上确界直接控制。一般指数的乘子问题还需另作研究，而全部 \(L^p\) 函数的变换将在温和分布中统一定义。
\end{chaptersummary}

\BookExercises

本组题目以 Plancherel 为计算工具，分别考察尺度、频带、乘子、算子收敛及空间与频率的同时集中。\Cref{b2:ex:04-four-sectors}研究 Fourier 算子的四个特征空间，为下一章 Hermite 函数作为 Fourier 特征函数的讨论提供算子观点。\Cref{b2:ex:04-support-volume}先建立空间投影与频率投影的复合估计，\cref{b2:ex:04-approximate-support}再将它用于有误差的集中；最后的\cref{b2:ex:04-interpolation-sign-moments}直接训练 Hausdorff--Young 延拓、尺度关系与符号平均。

% 自拟：Gauss范数检查非最优Hausdorff--Young常数及尺度不变性。
% 预留题头及首个公式所需空间，避免超链接锚点留在前页而题头移至后页。
\Needspace{6\baselineskip}
\begin{exercise}\label{b2:ex:04-gaussian-ratio}
设 \(d\) 为正整数，\(1<p<2\)、\(q=p/(p-1)\)，并对 \(a>0\) 定义实值函数 \(G_a(x)=e^{-\pi a|x|^2}\)（\(x\in\mathbb R^d\)）。计算 \(\|G_a\|_{L^p(\mathbb R^d)}\)、\(\|\widehat G_a\|_{L^q(\mathbb R^d)}\) 和二者的比值，说明为什么只改变 \(a\) 不能改善这个比值。该计算是否已经证明最佳 Hausdorff--Young 常数？
\end{exercise}
\begin{solution}
尺度公式给出
\[
 \|G_a\|_p=(ap)^{-d/(2p)},\quad
 \widehat G_a(\xi)=a^{-d/2}e^{-\pi|\xi|^2/a},\quad
 \|\widehat G_a\|_q=a^{-d/2}(q/a)^{-d/(2q)}.
\]
因 \(1/p+1/q=1\)，比值为 \(p^{d/(2p)}q^{-d/(2q)}\)，与 \(a\) 无关。这个数是任何普遍常数的下界；只计算一族函数，不能证明所有函数都满足取此常数的上界，因而不能单凭此计算宣称最佳性。
\end{solution}

% 自拟：紧频带产生光滑代表和定量导数界。
\Needspace{6\baselineskip}
\begin{exercise}\label{b2:ex:04-bandlimited}
设 \(d\) 为正整数，\(0<R<\infty\)、\(f\in L^2(\mathbb R^d;\mathbb C)\)，且 Plancherel 变换 \(\widehat f=\mathcal F_2f\) 几乎处处支撑在 \(B(0,R)\) 内。证明 \(f\) 有代表 \(h\in C^\infty(\mathbb R^d;\mathbb C)\)，并对每个多重指标 \(\alpha\in\mathbb N_0^d\) 有
\[
 \|\partial^\alpha h\|_\infty
 \leq(2\pi R)^{|\alpha|}|B(0,R)|^{1/2}\|f\|_2.
\]
\end{exercise}
\begin{solution}
有限频带上的 \(\widehat f\) 可积，乘任意多项式后仍可积。正号积分因此给出一个可任意次逐项求导的连续函数 \(h\)，且其 \(L^2\) 等价类由反演等于 \(f\)。求导后用 Cauchy--Schwarz：
\[
 |\partial^\alpha h(x)|
 \leq(2\pi R)^{|\alpha|}\int_{B(0,R)}|\widehat f|
 \leq(2\pi R)^{|\alpha|}|B(0,R)|^{1/2}\|\widehat f\|_2.
\]
最后用 Plancherel。这里的 \(R\) 有限是保证全部频率矩可积的关键。
\end{solution}

% 自拟：L2反演确实超出绝对可积函数域。
\Needspace{6\baselineskip}
\begin{exercise}\label{b2:ex:04-sinc}
在 \(\mathbb R\) 上使用 Lebesgue 测度，令 \(g(\xi)=\mathbf1_{[-1/2,1/2]}(\xi)\in L^2(\mathbb R;\mathbb C)\)，\(f=\mathcal F_2^{-1}g\)。求 \(f\) 的连续代表，证明 \(f\in L^2(\mathbb R;\mathbb C)\setminus L^1(\mathbb R;\mathbb C)\)，并计算 \(\|f\|_2\)。
\end{exercise}
\begin{solution}
有限频带积分给出 \(f(x)=\sin(\pi x)/(\pi x)\)，原点补为 \(1\)。Plancherel 给出 \(\|f\|_2=\|g\|_2=1\)。在每个 \([n+1/6,n+5/6]\) 上，\(|\sin(\pi x)|\geq1/2\)，因此其绝对积分至少为常数乘 \(\sum_{n\geq1}(n+1)^{-1}\)，故不属于 \(L^1\)。这不妨碍其 \(L^2\) 变换等于区间示性函数。
\end{solution}

% 自拟：平移群强连续却不按算子范数连续。
\Needspace{6\baselineskip}
\begin{exercise}\label{b2:ex:04-translation-operator}
设 \(d\) 为正整数，在 \(L^2(\mathbb R^d;\mathbb C)\) 上令 \((\tau_hf)(x)=f(x-h)\)（\(h\in\mathbb R^d\)），\(I\) 为恒等算子。证明对每个 \(f\in L^2(\mathbb R^d;\mathbb C)\)，\(\|\tau_hf-f\|_2\to0\) 当 \(h\to0\)；但对每个 \(h\ne0\)，\(\|\tau_h-I\|_{L^2\to L^2}=2\)。
\end{exercise}
\begin{solution}
平移的变换公式按稠密性推广到 \(L^2\)。于是
\[
 \|\tau_hf-f\|_2^2
 =\int|e^{-2\pi\mathrm{i}h\cdot\xi}-1|^2|\widehat f(\xi)|^2\,\mathrm{d}\xi\longrightarrow0
\]
由控制收敛成立。乘子范数则等于 \(e^{-2\pi\mathrm{i}h\cdot\xi}-1\) 的本性上确界；\(h\ne0\) 时，可在正测度频率带上使相位任意接近 \(-1\)，故该本性上确界为 \(2\)。前一结论固定函数，后一结论允许随 \(h\) 改变测试函数。
\end{solution}

% 自拟：频率投影和正性、伴随的分类。
\Needspace{6\baselineskip}
\begin{exercise}\label{b2:ex:04-multiplier-classification}
设 \(d\) 为正整数，\(m\in L^\infty(\mathbb R^d;\mathbb C)\)，并在 \(L^2(\mathbb R^d;\mathbb C)\) 上令 \(T_mf=\mathcal F_2^{-1}(m\cdot\mathcal F_2f)\)。求 \(T_m^*\)，并分别刻画 \(T_m\) 自伴、非负、酉和正交投影的充要条件。这里非负指 \(\langle T_mf,f\rangle\) 对每个 \(f\in L^2(\mathbb R^d;\mathbb C)\) 都是非负实数。
\end{exercise}
\begin{solution}
在频率侧计算内积，得 \(T_m^*=T_{\overline m}\)。所以自伴当且仅当 \(m\) 几乎处处为实数。非负当且仅当 \(m\geq0\) 几乎处处，充分性由积分给出。为证必要性，若虚部不为零，可在某个有限正测度集 \(E\) 上使它同号，且其绝对值有正下界；以 \(\widehat f=\mathbf1_E\) 测试，便使 \(\langle T_mf,f\rangle\) 的虚部非零。因此符号先须几乎处处为实数。若实部在正测度集合上为负，同样取有限正测度的负值水平集作测试，就得到负的二次型，矛盾。又因 \(T_m^*T_m=T_{|m|^2}\)，酉性等价于 \(|m|=1\) 几乎处处，此时逆符号为 \(\overline m\)。正交投影等价于自伴且幂等，故 \(m\) 几乎处处只能取 \(0,1\)，即为某个可测集的示性函数。
\end{solution}

% 自拟：傅里叶算子的四个代数谱分量；无须预借Hermite函数。
\Needspace{6\baselineskip}
\begin{optionalexercise}\label{b2:ex:04-four-sectors}
设 \(d\) 为正整数，\(H=L^2(\mathbb R^d;\mathbb C)\)，\(\mathcal F_2\colon H\to H\) 为 Plancherel 变换，\(I\) 为 \(H\) 上的恒等算子。对 \(\lambda\in\{1,-1,\mathrm{i},-\mathrm{i}\}\)，定义
\[
 Q_\lambda=\dfrac14\sum_{j=0}^3\lambda^{-j}\mathcal F_2^j.
\]
其中 \(\mathcal F_2^0=I\)。证明这些算子是两两正交的正交投影，其和为 \(I\)，且其值域等于 \(\{\,f\in H\mid\mathcal F_2f=\lambda f\,\}\)。说明 \(\lambda=\pm1\) 的值域只含几乎处处意义下的偶函数，\(\lambda=\pm\mathrm{i}\) 的值域只含几乎处处意义下的奇函数。
\end{optionalexercise}
\begin{solution}
利用 \(\mathcal F_2^4=I\) 移动求和指标，可得 \(\mathcal F_2Q_\lambda=\lambda Q_\lambda\)。若 \(\mathcal F_2f=\mu f\)，则 \(Q_\lambda f=\frac14\sum_{j=0}^3(\mu/\lambda)^jf\)，当 \(\mu=\lambda\) 时等于 \(f\)，否则为零。并且四个 \(Q_\lambda\) 的和为 \(I\)，因为四次单位根的非零次幂和为零。这已给出幂等性、相互乘积为零和值域描述。酉性使不同单位模特征值的特征向量正交：\(\langle f,g\rangle=\lambda\overline\mu\langle f,g\rangle\)。故这些投影确为正交投影。最后 \(Rf=\mathcal F_2^2f=\lambda^2f\)，得到奇偶性。
\end{solution}

% 自拟：空间投影与频率投影的复合范数估计，作为近似支集题的算子预备。
\Needspace{6\baselineskip}
\begin{exercise}\label{b2:ex:04-support-volume}
设 \(d\) 为正整数，\(A,B\subset\mathbb R^d\) 是 Lebesgue 可测集，且 \(|A|,|B|<\infty\)。在 \(L^2(\mathbb R^d;\mathbb C)\) 上定义
\[
 P_Af=\mathbf1_A\cdot f,
 \quad Q_Bf=\mathcal F_2^{-1}(\mathbf1_B\cdot\mathcal F_2f).
\]
证明 \(P_A,Q_B\) 是正交投影，并证明
\[
 \|Q_BP_A\|_{L^2\to L^2}\leq|A|^{1/2}\cdot|B|^{1/2}.
\]
说明：如果某个非零 \(f\in L^2(\mathbb R^d;\mathbb C)\) 同时满足 \(P_Af=f\)、\(Q_Bf=f\)，那么这一初等估计只能给出必要条件 \(|A|\cdot|B|\geq1\)，并未构造这样的函数。
\end{exercise}
\begin{solution}
乘以示性函数 \(\mathbf1_A\) 的算子自伴且幂等，所以 \(P_A\) 是正交投影。频率侧乘以 \(\mathbf1_B\) 同样是正交投影；与酉算子 \(\mathcal F_2\) 共轭便得到 \(Q_B\)。因此两者的算子范数都不超过 \(1\)。对任意 \(u\in L^2(\mathbb R^d;\mathbb C)\)，Cauchy--Schwarz 给出
\[
 \|P_Au\|_1\leq|A|^{1/2}\cdot\|P_Au\|_2
 \leq|A|^{1/2}\cdot\|u\|_2.
\]
所以 \(P_Au\in L^1\cap L^2\)，其普通 Fourier 变换与 Plancherel 变换几乎处处相等。由酉性及 \(L^1\to L^\infty\) 估计，
\[
 \begin{aligned}
 \|Q_BP_Au\|_2
 &=\|\mathbf1_B\cdot\widehat{P_Au}\|_2\\
 &\leq|B|^{1/2}\cdot\|\widehat{P_Au}\|_\infty\\
 &\leq|B|^{1/2}\cdot\|P_Au\|_1
 \leq|A|^{1/2}\cdot|B|^{1/2}\cdot\|u\|_2.
 \end{aligned}
\]
对全部单位向量取上确界即得算子范数估计；\(|A|=0\) 或 \(|B|=0\) 的情形也包含在内。若另有题目末尾的非零 \(f\)，则 \(Q_BP_Af=f\)，上式成为 \(\|f\|_2\leq|A|^{1/2}\cdot|B|^{1/2}\cdot\|f\|_2\)，除以正数 \(\|f\|_2\) 即得所述条件。
\end{solution}

作为书外补充，Benedicks 的《The support of functions and distributions with a spectral gap》\cite[p. 286]{Benedicks1984SpectralGap}中 Theorem 1 给出更深的结论：对任意正整数 \(d\)，若 \(f\in L^2(\mathbb R^d;\mathbb C)\) 及 \(\mathcal F_2f\) 分别几乎处处支撑于有限测度可测集，则 \(f=0\) 几乎处处。该定理的证明见所引论文。

% 自拟：把空间、频率投影估计用于有误差的同时集中，假设可由实际函数实现。
\Needspace{6\baselineskip}
\begin{optionalexercise}\label{b2:ex:04-approximate-support}
设 \(d\) 为正整数，\(f\in L^2(\mathbb R^d;\mathbb C)\)、\(\|f\|_2=1\)，\(\widehat f=\mathcal F_2f\)，\(A,B\subset\mathbb R^d\) 是有限 Lebesgue 测度的可测集。设 \(\varepsilon,\delta\geq0\)、\(\varepsilon+\delta<1\)，且
\[
 \|f\cdot\mathbf1_{A^c}\|_2\leq\varepsilon,
 \quad\|\widehat f\cdot\mathbf1_{B^c}\|_2\leq\delta.
\]
证明 \(|A|^{1/2}\cdot|B|^{1/2}\geq1-\varepsilon-\delta\)。
\end{optionalexercise}
\begin{solution}
令 \(P_Af=\mathbf1_A\cdot f\)、\(Q_Bf=\mathcal F_2^{-1}(\mathbf1_B\cdot\widehat f)\)。两者均为正交投影，故
\[
 \|f-Q_BP_Af\|_2
 \leq\|f-Q_Bf\|_2+\|Q_B(f-P_Af)\|_2\leq\delta+\varepsilon.
\]
\Cref{b2:ex:04-support-volume}已经对任意输入证明了复合算子范数的上界。对当前单位向量使用它，得到
\[
 \|Q_BP_Af\|_2
 \leq\|Q_BP_A\|_{L^2\to L^2}\cdot\|f\|_2
 \leq|A|^{1/2}\cdot|B|^{1/2}.
\]
由三角不等式，\(1\leq\varepsilon+\delta+\|Q_BP_Af\|_2\)，得到所求。

这些条件确能由非零函数实现。例如取归一化的非零 \(f\in C_c^\infty(\mathbb R^d;\mathbb C)\)，以一个有界集 \(A\) 包含其支集，空间误差就是零。又因 \(\widehat f\in L^2\)，可取足够大的有限半径球 \(B\)，使频率侧尾部的平方积分任意小。这说明近似集中允许非零对象，而无需要求两个精确支集都有限。
\end{solution}

% 自拟：具体非Gauss尾函数的Hausdorff--Young延拓与尺度；加权符号和的精确四矩及矩插值，不复现正文波包反例。
\Needspace{6\baselineskip}
\begin{exercise}\label{b2:ex:04-interpolation-sign-moments}
\begin{enumerate}
\item\label{b2:item:04-hy-nongaussian-scaling} 在一维令 \(h(x)=(1+|x|)^{-1}\)，\(1<p<2\)、\(q=p/(p-1)\)。证明 \(h\in L^p(\mathbb R;\mathbb C)\setminus L^1(\mathbb R;\mathbb C)\)。记 \(H=\mathcal F_ph\in L^q(\mathbb R;\mathbb C)\) 为 Hausdorff--Young 延拓，证明
\[
 \|H\|_q\leq\left(\frac2{p-1}\right)^{1/p}.
\]
对 \(a>0\) 令 \(h_a(x)=h(ax)\)，在 \(L^q\) 意义下证明 \(\mathcal F_ph_a(\xi)=a^{-1}H(\xi/a)\)，并核对输入、输出范数具有同一个尺度因子 \(a^{-1/p}\)。不能把 \(h\) 直接代入绝对收敛 Fourier 积分。
\item\label{b2:item:04-weighted-sign-moments} （选做）设 \(N,d\) 为正整数，\(c_1,\ldots,c_N\in\mathbb C\) 不全为零，\(S=\sum_{j=1}^N|c_j|^2\)。在全部 \(2^N\) 组 \(\epsilon=(\epsilon_1,\ldots,\epsilon_N)\in\{-1,1\}^N\) 上取等权平均 \(\mathbb E\)，令 \(Z=\sum_j\epsilon_jc_j\)。求 \(\mathbb E|Z|^2\)、\(\mathbb E|Z|^4\)，并用 Hölder 不等式证明
\[
 \mathbb E|Z|\geq\sqrt{S/3}.
\]
进一步，给定任意 \(t_1,\ldots,t_N\in\mathbb R^d\) 及非负 \(w\in L^1(\mathbb R^d;\mathbb R)\)、\(\int_{\mathbb R^d}w>0\)，证明存在同一组固定符号 \(\sigma_j\in\{-1,1\}\)，使
\[
 \int_{\mathbb R^d}w(\xi)\cdot
 \left|\sum_{j=1}^N\sigma_jc_je^{-2\pi\mathrm i t_j\cdot\xi}\right|\,\mathrm d\xi
 \geq\sqrt{S/3}\int_{\mathbb R^d}w.
\]
说明为何不能对每个频率分别选一组符号后直接积分，并解释 \(c_j=1\) 时 \(N^{1/2}\) 增长的来源。

最后取 \(2<r\leq\infty\)、\(1/r+1/s=1\)，以及非零 \(\varphi\in C_c^\infty(\mathbb R^d;\mathbb C)\)。选 \(L>0\) 使 \(\varphi(\cdot-jLe_1)\) 的支集互不相交。将上述结论用于 \(c_j=1\)、\(t_j=jLe_1\)、\(w=|\widehat\varphi|^s\)，证明可以选同一组符号使
\[
 \begin{gathered}
 f_N(x)=\sum_{j=1}^N\sigma_j\varphi(x-jLe_1),\\
 \|f_N\|_r=N^{1/r}\|\varphi\|_r,\qquad
 \|\widehat {f_N}\|_s\geq\sqrt{N/3}\,\|\widehat\varphi\|_s.
 \end{gathered}
\]
其中 \(N^{1/\infty}=1\)。提示：在有限测度 \(|\widehat\varphi|^s\,\mathrm d\xi\) 上使用 Hölder 不等式。
\end{enumerate}
\end{exercise}
\begin{solution}
第一问，直接积分给出
\[
 \|h\|_p^p=2\int_0^\infty(1+x)^{-p}\,\mathrm dx=\frac2{p-1}.
\]
将指数换成 \(1\) 则积分发散。因此所用变换必须按\cref{b2:cor:04-lp-extension}的 \(L^p\) 延拓理解，该推论立即给出所求范数界。为证明尺度公式，令 \(h_n=h\cdot\mathbf1_{[-n,n]}\)，则 \(h_n\in L^1\cap L^p\)、\(h_n\to h\) 于 \(L^p\)，并且 \(\widehat{h_n}\to H\) 于 \(L^q\)。对这些可积截断先用伸缩公式：
\[
 \widehat{h_n(a\,\cdot)}(\xi)=a^{-1}\widehat{h_n}(\xi/a).
\]
左侧因输入在 \(L^p\) 中趋于 \(h_a\)，故在 \(L^q\) 中趋于 \(\mathcal F_ph_a\)；右侧在 \(L^q\) 中趋于 \(a^{-1}H(\cdot/a)\)。这里每个固定 \(a>0\) 的伸缩都连续，因为变量代换给出相应范数的精确尺度因子。极限唯一性便证明了公式，并得到
\[
 \|h_a\|_p=a^{-1/p}\left(\frac2{p-1}\right)^{1/p},
 \quad
 \|\mathcal F_ph_a\|_q
 =a^{-1+1/q}\|H\|_q=a^{-1/p}\|H\|_q.
\]

第二问，符号平均使含某个 \(\epsilon_j\) 奇数次的单项全部消失，所以
\[
 \mathbb E|Z|^2=S,
 \quad
 \mathbb E|Z|^4
 =2S^2+\left|\sum_jc_j^2\right|^2-2\sum_j|c_j|^4
 \leq3S^2.
\]
四阶等式可以按三种成对划分核验：两种给出 \(S^2\)，另一种给出 \(|\sum_jc_j^2|^2\)；四个指标全相等的项被计了三次，故减去两次 \(\sum_j|c_j|^4\)。最后用 \(|\sum_jc_j^2|\leq\sum_j|c_j|^2=S\) 得到上界。对有限符号空间使用指数 \(3/2\) 与 \(3\) 的 Hölder 不等式，
\[
 S=\mathbb E\bigl(|Z|^{2/3}\cdot|Z|^{4/3}\bigr)
 \leq(\mathbb E|Z|)^{2/3}(\mathbb E|Z|^4)^{1/3}.
\]
因 \(S>0\)，整理即得 \(\mathbb E|Z|\geq S^{3/2}/(\mathbb E|Z|^4)^{1/2}\geq\sqrt{S/3}\)。这一步把二、四矩的控制转成一阶绝对矩的下界，无须重复正文的阈值事件论证。

对每个固定 \(\xi\)，把系数换成 \(c_je^{-2\pi\mathrm i t_j\cdot\xi}\)，它们的模不变，所以同一个 \(\sqrt{S/3}\) 下界仍适用。乘以 \(w(\xi)\) 后积分，再交换有限平均与积分，得到
\[
 2^{-N}\sum_{\epsilon\in\{-1,1\}^N}
 \int_{\mathbb R^d}w(\xi)\cdot
 \left|\sum_j\epsilon_jc_je^{-2\pi\mathrm i t_j\cdot\xi}\right|\,\mathrm d\xi
 \geq\sqrt{S/3}\int_{\mathbb R^d}w.
\]
每个积分都不超过 \((\sum_j|c_j|)\int w<\infty\)。至少一个被平均的积分不小于平均值，选择那一组符号即得结论；这组符号在整个积分中保持固定。逐点另选符号可能随 \(\xi\) 改变，不能保证来自同一个有限和。等系数时 \(S=N\)，平方根增长正是二矩可加而四矩至多与其平方同阶的结果。

对最后的应用，光滑紧支集函数的变换可积且有界，所以 \(w=|\widehat\varphi|^s\in L^1\)；Fourier 唯一性保证 \(W\coloneqq\int w>0\)。支集分离给出 \(\|f_N\|_r=N^{1/r}\|\varphi\|_r\)，包括 \(r=\infty\)。令 \(Z_\sigma(\xi)=\sum_j\sigma_je^{-2\pi\mathrm i jL\xi_1}\)，所选符号满足 \(\int w|Z_\sigma|\geq\sqrt{N/3}\,W\)。当 \(s>1\) 时，Hölder 不等式给出
\[
 \int w|Z_\sigma|
 \leq\left(\int w|Z_\sigma|^s\right)^{1/s}W^{1-1/s};
\]
\(s=1\) 时同一式就是等式。因此
\[
 \|\widehat {f_N}\|_s
 =\left(\int w|Z_\sigma|^s\right)^{1/s}
 \geq\sqrt{N/3}\,W^{1/s}
 =\sqrt{N/3}\,\|\widehat\varphi\|_s.
\]
这由加权符号平均重新得到正文的增长机制，使用的是矩插值而非阈值事件；与输入的 \(N^{1/r}\) 比较，就看出 \(r>2\) 时范数界为何失效。
\end{solution}
